\section{Arquitectura funcional M1--M8}

La nomenclatura M1--M8 se conserva como referencia arquitectónica procedente del modelo histórico. No todos los módulos se encuentran actualmente implementados con el mismo grado de madurez en Python.

\begin{table}[H]
\centering
\caption{Funciones y estado de los módulos M1--M8.}
\begin{tabular}{p{0.15\textwidth}p{0.43\textwidth}p{0.30\textwidth}}
\toprule
Módulo & Función & Estado actual \\
\midrule
M1 & Instrumentación y observabilidad & Reconstrucción funcional \\
M2 & Convertidor de potencia & Simplificado en la reconstrucción \\
M3 & Interfaz eléctrica y transformaciones & Reconstrucción funcional \\
M4 & Planta PMSM & Modelo base establecido \\
M5 & Evidencia, métricas y validación & Consistencia/reproducibilidad; validación física pendiente \\
M6 & Control FOC & Funcional/reconstruido \\
M7 & Inteligencia, adaptación y predicción & En desarrollo \\
M8 & Supervisión superior y seguridad & Preparado; implementación final pendiente \\
\bottomrule
\end{tabular}
\end{table}

El flujo principal del sistema es:
\begin{equation}
\text{M6 FOC}\rightarrow\text{M3}\rightarrow\text{M2}\rightarrow\text{M4 PMSM}\rightarrow\text{retroalimentación}.
\end{equation}

La capa M7 deberá interactuar mediante contratos de frontera y no sustituirá al FOC como controlador principal.
